\documentclass[aspectratio=169,  notheorems, sOuRcEs]{RUCPresentation}

%\usepackage{pgfforeach}
% \usepackage{foreach}

\newcounter{chapter}


\usepackage{tikz}
\usetikzlibrary{calc}
\usetikzlibrary{angles}
\usetikzlibrary{positioning}

\usetikzlibrary{arrows.meta}
\usetikzlibrary{shapes.geometric}
\usetikzlibrary{overlay-beamer-styles}

\ProvideDocumentCommand\pgfmathsetnewmacro{ m m }{ \newcommand*{#1}{} \pgfmathsetmacro{#1}{#2} }
\providecommand*{\pgfmathsetnewlengthmacro}[2]{ \newcommand*{#1}{} \pgfmathsetlengthmacro{#1}{#2} }

\tikzset{>=Stealth} % for good looking LaTeX arrow head
% \tikzset{overlay, remember picture}


\usepackage{myMathQA}
\usepackage{graphicx}
\usepackage{siunitx}

\usepackage{hyperref}

% \usepackage[style=authoryear]{biblatex}
\usepackage{biblatex}
\addbibresource{mscbib.bib}

\DeclareSourcemap{
    \maps[datatype=bibtex]{
        \map[overwrite]{
            \step[fieldset=note, null]
        }
        \map[overwrite]{
            \step[fieldset=ppn_gvk, null]
        }
        \map[overwrite]{
            \step[fieldset=file, null]
        }
        \map[overwrite]{
            \step[fieldset=doi, null]
        }
        \map[overwrite]{
            \step[fieldset=issn, null]
        }
        \map[overwrite]{
            \step[fieldset=isbn, null]
        }
        \map[overwrite]{
            \step[fieldset=pagetotal, null]
        }
    }
}

% \DeclareCiteCommand{\beamerfootcite}
%   {}
%   {\footnote[frame]{\printbibliography[keys=\thefield{entrykey}]}}
%   {\multicitedelim}
%   {\usebibmacro{postnote}}

\ProvideDocumentCommand{\beamerfootcite}{ m }{\footnote[frame]{\ \fullcite{#1}}}
\renewcommand*{\nameyeardelim}{\addcomma\space} % cleaner authoryear style

% \usepackage{animate}



\def\k{\raisebox{1ex}{\rotatebox{180}{k}}}
% \newsavebox\tempboxbeamer

\usetheme{RUCTHEME}
% \usetheme{Warsaw}

\makeatletter
\renewcommand\@makefnmark{\hbox{\@textsuperscript{\usebeamercolor[fg]{footnote mark}\usebeamerfont*{footnote mark}[\@thefnmark]}}}
% \renewcommand\@makefntext[1]{{\usebeamercolor[fg]{footnote mark}\usebeamerfont*{footnote mark}[\@thefnmark]}\usebeamerfont*{footnote} #1}
\makeatother

\hypersetup{colorlinks = false, citecolor=TextColor}
\AtBeginBibliography{\small\color{TextColor}}


\def\dumbmatrixtwobytwo[#1,#2]{%
    \count6=#1\relax \count7=#2\relax%
    \count2=#1\relax \count3=#2\relax \count4=#1\relax \count5=#2\relax%
    \multiply\count2 by \count6\relax \multiply\count4 by \count7\relax%
    \multiply\count3 by \count6\relax \multiply\count5 by \count7\relax%
    \ensuremath{%
        \left[%
            \begin{matrix}%
                \number\count2 & \number\count4 \\%
                \number\count3 & \number\count5%
            \end{matrix}%
        \right]%
    }
}


\def\dumbmatrixthreebythree[#1,#2,#3]{%
    \count10=#1\relax \count11=#2\relax \count12=#3\relax
    \count1=#1\relax \count4=#1\relax \count7=#1\relax%
    \count2=#2\relax \count5=#2\relax \count8=#2\relax%
    \count3=#3\relax \count6=#3\relax \count9=#3\relax%
    \multiply\count1\count10\relax \multiply\count4\count11\relax \multiply\count7\count12\relax%%
    \multiply\count2\count10\relax \multiply\count5\count11\relax \multiply\count8\count12\relax%%
    \multiply\count3\count10\relax \multiply\count6\count11\relax \multiply\count9\count12\relax%%
    \ensuremath{%
        \left[%
            \begin{matrix}%
                \number\count1 & \number\count4 & \number\count7 \\%
                \number\count2 & \number\count5 & \number\count8 \\%
                \number\count3 & \number\count6 & \number\count9%
            \end{matrix}%
        \right]%
    }
}

\usepackage{multirow}
\newcounter{reaction}
\ProvideDocumentCommand\reactionterm{ o m m m m }{% --> Function for reaction lines
    \stepcounter{reaction} %
    \IfNoValueTF{#1}
    { \arabic{reaction} & #3 & #4 & \(t^+_{\arabic{reaction}}\of{\vec{x}}\) & #2 & \(\br{#5}\) \\ \hline }
    { \multirow{2}*{\arabic{reaction}} & #3 & #4 & \(t^+_{\arabic{reaction}}\of{\vec{x}}\) & #2 &  \multirow{2}*{\(\br{#5}\)} \\
                                       & #4 & #3 & \(t^-_{\arabic{reaction}}\of{\vec{x}}\) & #1 &                             \\ \hline }
}

\ProvideDocumentEnvironment{reactiontable}{}%
{ \begin{center} \begin{tabular}{c|r@{\ \(\to\)\ }lr@{\ \(=\)\ }lc} %
    \(I\) & \multicolumn{2}{c}{`Reaction'} %
          & \multicolumn{2}{c}{`Rate Law'} %
          & \(\vec{r}_I^{\mathrm{T}}\) \\\hline }%
{ \end{tabular} \end{center} \setcounter{reaction}{0} }


% Defining a new coordinate system for the page:
%
% --------------------------
% |(-1,1)    (0,1)    (1,1)|
% |                        |
% |(-1,0)    (0,0)    (1,0)|
% |                        |
% |(-1,-1)   (0,-1)  (1,-1)|
% --------------------------

% For beamer title
\newcommand\makebeamertitle{\frame{\maketitle}}%

% For sections?
\graphicspath{{./images/}{../Thesis/images/phase/}{../Thesis/images/ode/}{../Thesis/images/fpe/}}


\def\cfp#1{c_#1^*}


\title{%
    An exploration on the differences of stochastic and
    deterministic models
}
% \subtitle{  Mathematical\ \ Bioscience a MSc Thesis }
\subtitle{  My MSc Thesis Summerized }
% \newline supervisor: Morten Anderson}

\author{by Anakin}

%\mode<presentation>


\begin{document}


\makebeamertitle


\section{Introduction}

\begin{frame}{The Research Question}

    \begin{itemize}
        \item So why am \emph{I} interested in this
    \end{itemize}
    \pause
    \begin{enumerate}[<+->]
        \item Nothing noble
        \item I wanted to find an excuse to learn about stochasticity
        \item I believe there is an inherent relationship between
            statistical manifolds and genetic phenotypes
    \end{enumerate}

    \pause
    \begin{block}{As a consequence, my research question}
        How can stochastic gene copy number dynamics underlying bacterial
        heteroresistance be accurately captured using a Fokker-Planck formulation,
        and what insights does this stochastic framework provide beyond deterministic ODE models.
    \end{block}

\end{frame}

% Mechanisms of heteroresistance:
%     Caused by frequent mutations other than spontaneous gene amplifications:
%         Point mutations
%         Insertion sequence insertions
%         Small deletions
%     Resistance mutations with low or no fitness cost:
%         Stable even in the absence of antibiotics
%         Leads to stable heteroresistance
%     High fitness cost associated with resistance mutation:
%         Resistant phenotype can revert in the absence of antibiotics
%         Second-site compensatory mutations can increase fitness
%         Concomitant loss of resistance
%         Generates a mixed population:
%             Rare resistant cells with reduced fitness (contain a resistance mutation)
%             Main population of susceptible cells with higher fitness (contain a resistance mutation and a compensatory mutation)
%         Revertants often have different genotypes and/or phenotypes from the original isolate
%     Unstable resistance:
%         Caused by intrinsic instability of the resistance mutation itself:
%             Spontaneous unstable gene amplifications
%             Causes unstable heteroresistance due to the cost and high instability of gene amplifications
%         Revertants have the same genotype and phenotype as the original isolate


\section{Motivation}

\subsection{Heteroresistence}

\begin{frame}{Heteroresistence}

    % Text\beamerfootcite{Levin2024} supposed to be citation

    % The presence of resistance genes (mutated or acquired genes) in bacterial
    % cells typically leads to measurable antibiotic resistance, creating a
    % generally predictable relationship between a bacterium's genetic makeup
    % (genotype) and its observed antibiotic susceptibility (phenotype).
    % However, there are important exceptions.
    % Sometimes even genetically identical bacterial cells within the same
    % population can display phenotypic heterogeneity (variability) in their
    % responses to antibiotics, complicating the clear-cut classification into
    % susceptible or resistant categories. Heteroresistance is a notable example
    % of this variability, wherein a bacterial isolate predominantly susceptible
    % to an antibiotic includes smaller subpopulations of cells with reduced
    % antibiotic susceptibility.


    \begin{itemize}[<+->]
        \item Presence of resistance genes in bacterial cells leads
            to measurable antibiotic resistance
            \begin{itemize}[<.->]
                \item Generally predictable relationship between bacterial
                    genotype and antibiotic resistance
                \item Exceptions exist
            \end{itemize}
        \item Genetically identical bacterial cells can display phenotypic
            heterogeneity in responses to antibiotics.
            % \begin{itemize}
            %     \item Complicates classification into susceptible or
            %         resistant categories
            % \end{itemize}
        \item Heteroresistance: a bacterial isolate predominantly
            susceptible to an antibiotic includes smaller subpopulations
            of cells with reduced antibiotic susceptibility\beamerfootcite{Andersson2019}
    \end{itemize}

\end{frame}

\begin{frame}

    What does heteroresistance look like?
    \begin{columns}[T]
        \column{0.5\framewidth}

        \begin{itemize}
            \item Resistance mutations with low or no fitness cost
                \begin{itemize}
                    \item Insertions
                    \item Deletions
                \end{itemize}
        \end{itemize}

        \column{0.5\framewidth}

        \begin{itemize}
            \item High fitness cost associated with resistance mutation
                \begin{itemize}
                    \item Insertions
                    \item Deletions
                \end{itemize}
        \end{itemize}

    \end{columns}

\end{frame}

\subsection{Copy Numbers}

\begin{frame}{Gene Copy Numbers}
\end{frame}


\begin{frame}{Fitness Doesn't Matter?}



\end{frame}



\section{What do we know}


\begin{frame}{ODE system}

    \tikz[overlay] \node<1-> at (frame cs:-0.65,-1.35) {
        \begin{tikzpicture}[ scale = 1,
            draw = black!80,
            thick,
            shorten >=2pt,
            overlay
            ]

            \node[draw=black, circle] (c1) at (-1, 0) {\(c_1\)};
            \node[draw=black, circle] (c2) at ( 1, 0) {\(c_2\)};

            \node[draw=black] (m) at ( -1, -1.20) {\(m\)};
            \node[] (source) at ( -2, -1.20) {\(m_0\)};

            \draw[->] (c2) [bend left]to node[midway, below] {\(w_2\)} (c1);
            \draw[->] (c1) [bend left]to node[midway, above] {\(w_1\)} (c2);

            \draw[->] (c1)  [loop left]to node[midway, left]  {\(k_1\)} (c1);
            \draw[->] (c2) [loop right]to node[midway, right] {\(k_2\)} (c2);

            \draw[-|] (m) to node[pos=0.35] (mm) {} (c1);
            \draw[-|, rounded corners] (c2) |- node[midway,anchor=south west] {\(q_m\)} (m);
            % \draw[-|] (c2) to[out=270, in=0, looseness=2] node[midway,anchor=north west] {\(q_m\)} (m);

            \node[anchor=east] at (mm) {\(q_1\)};

            \draw[->] (source) -- (m);

    \end{tikzpicture}};


    \begin{columns}
        \column{0.5\framewidth}
        \vspace{8em}\par

        \begin{align*}
            \ode{c_1} &= k_1 c_1\left({n_1 - c_t \over n_1}\right) - w_1 c_1 + w_2 c_2 - q_1 m c_1 \nonumber \\
            \ode{c_2} &= k_2 c_2\left({n_2 - c_t \over n_2}\right) - w_2 c_2 + w_1 c_1\nonumber\\
            \ode{m}   &= m_0 - q_m m c_2 \nonumber
        \end{align*}

        \column{0.5\framewidth}

    \end{columns}



    \begin{tikzpicture}[overlay, remember picture]

        \node<2-> at (frame cs:0.5, 0.34) {
                \includegraphics[width=0.4\textwidth]
                {ode:jesper:I6:C1:T17481880809603815.pdf}
            };
        \node<2-> at (frame cs:0.5, -0.40) {
                \includegraphics[width=0.4\textwidth]
                {ode:jesper:I32:C1:T17481997939903586.pdf}
            };

    \end{tikzpicture}

\end{frame}

\begin{frame}{Analytical solution}


    \begin{columns}[T]
        \column{0.5\framewidth}
        \begin{itemize}
            \alt<1>{
                \item Skipping the details
                }{
                \item Soure: trust me
                }
            \item We recognize that
        \end{itemize}


        \column{0.5\framewidth}


    \end{columns}

    \begin{tikzpicture}[overlay, remember picture]

        \node<2-> at (frame cs:0.5, 0.34) {
                \includegraphics[width=0.4\textwidth]
                {ode:jesper:I6:C1fp:T17481880809603815.pdf}
            };
        \node<2-> at (frame cs:0.5, -0.40) {
                \includegraphics[width=0.4\textwidth]
                {ode:jesper:I32:C1fp:T17481997939903586.pdf}
            };

    \end{tikzpicture}

\end{frame}




\section{Stochasticity}

\begin{frame}{Stochasticity}

    \begin{itemize}
        \item Stochasticity is a touchy topic mathematically
            \begin{itemize}[<.->]
                \item
            \end{itemize}
    \end{itemize}

    So what some \emph{lunatic} did, is glue a Wiener process to the end of our
    differential.
    With some multivariate function \(B\), acting on \(W_t\) as a scaling factor
    for the variance.
    By substitution of variables, we can redefine the sides of the derivative,
    \begin{align}
        %\ode x = A\of{x,t} + B\of{x,t}
        \ode x = A\of{x\of t, t} \implies \dd{x} &= A\of{x,t} \dd{t}
        \shortintertext{Throw our randomness in there}
        \dd{x} &= A\of{x,t} \dd{t} + \sqrt{  B\of{x,t} } \ \dd{W_t}
    \end{align}
    from allowing our ODE to be influenced by stochastic noise, we now
    have a SDE.

    More formally known as the the \emph{Langevin equation},

\end{frame}

\begin{frame}{The Master Equation}


    The dynamics of a probability distribution with the Markov property
    have a \emph{master equation}\beamerfootcite{Gardiner1987}
    \begin{equation} \label{eq.mastereq}
        \pde* P\of{x,t} =\int_{-\infty}^{\infty}
        \left[W_t\of{x|x'}P\of{x',t} - W_t\of{x'|x} P\of{x,t}\right]
        \dd{x'}
    \end{equation}
    \begin{itemize}
        \item  \(W_t\of{x|x'}\) is the rate of state transitions from \(x'\)
        \item  \(P\of{x, t}\) is the probability of \(x\) at the given time
    \end{itemize}

\end{frame}

\begin{frame}{The Chemical Master Equation}


    \setcounter{footnote}{1}
    But in the non-continous form, we can make a \emph{chemical}
    master equation\beamerfootcite{Gardiner1987}
    \begin{equation} \label{eq.chemmastereq}
        \pde* P\of{x,t}
        =
        \sum_{j}^{}\br{
            \left[a_j\of{x|x'} P\of{x',t} - a_j\of{x'|x} P\of{x,t}\right]
        }
    \end{equation}

    \begin{itemize}
        \item \(a_j\) represents a `propensity'
        \item \(\vec{v}_j\) a state change vector
        \item \(\vec{n}_j\) a state change vector
        \item Basically just a discrete version of the master equation
            \begin{itemize}
                \item That was a lie
            \end{itemize}
    \end{itemize}

\end{frame}

\begin{frame}{Fokker-Planck}

\end{frame}



\subsection{The Ordinary Differential System}

\begin{frame}{Chemicafying it}


    \begin{columns}[T]
        \column{0.3\textwidth}
        \begin{itemize}
            \item Recontextualizing the compartment model into kinetics
        \end{itemize}
        \begin{align*}
            c_1 &\xrightleftharpoons[w_2]{w_1} c_2 \\
            c_1 + N &\xrightleftharpoons[k_{-1}]{k_1} 2c_1\\
            c_2 + N &\xrightleftharpoons[k_{-2}]{k_2} 2c_2
        \end{align*}

        \column{0.7\textwidth}
        \begin{itemize}
            \item Letting  \(\vec{r}_1 = \br{-1,1}^{\mathrm{T}} \) be our transition vector,
            \item The first kinetic reaction \(c_1 \xrightleftharpoons[w_2]{w_1} c_2 \)
            becomes
        \end{itemize}
        \begin{align*}
                T\of{ \vec{c} \to \vec{c} + \vec{r}_1} &= w_1 {c_1!\over {c_1 - 1}!}{c_2!\over {c_2 - 0}!} \\
            &\leadsto t_1^+\of{\vec{c}} = w_1 c_1 \\
                T\of{ \vec{c} \to \vec{c} - \vec{r}_1} &= w_2 {c_1!\over {c_1 - 0}!}{c_2!\over {c_2 - 1}!} \\
            &\leadsto t_1^-\of{\vec{c}} = w_2 c_2
        \end{align*}

    \end{columns}

\end{frame}

\begin{frame}

    \begin{columns}[T]
        \column{0.5\textwidth}

        \begin{itemize}
            \item The same process for the rest of the terms leads to a table
        \end{itemize}

        \begin{reactiontable}
            \reactionterm[\(w_2 c_2\)]{\(w_1 c_1\)}{\(c_1\)}{\(c_2\)}{-1, 1}
            \reactionterm[\(k_{-1} c_1^2\)]{\(k_1 n c_1\)}{\(c_1+N\)}{\(2 c_1\)}{ 1, 0}
            \reactionterm[\(k_{-2} c_2^2\)]{\(k_2 n c_2\)}{\(c_2+N\)}{\(2 c_2\)}{ 0, 1}
        \end{reactiontable}

        \column{0.5\textwidth}

    \end{columns}


    \begin{tikzpicture}[overlay, remember picture]


        \node at (frame cs:0.55, 0.45) {
                \includegraphics[width=0.38\framewidth, page=1]
                {ode/ssa:ode-3-2:I19:C2fp:T1747486672.pdf}
            };

        \node at (frame cs:0.55,-0.35) {
                \includegraphics[width=0.38\framewidth, page=3]
                {ode/ssa:ode-3-2:I19:C2fp:T1747486672.pdf}
            };

    \end{tikzpicture}

\end{frame}

\begin{frame}{We add Smothering}

    \begin{columns}[T]
        \column{0.5\textwidth}

        \begin{itemize}
            \item With smothering
        \end{itemize}

        \alt<1>{
            \begin{reactiontable}
                \reactionterm[\(w_2 c_2\)]{\(w_1 c_1\)}{\(c_1\)}{\(c_2\)}{-1, 1}
                \reactionterm[\(k_{-1} c_1^2\)]{\(k_1 n c_1\)}{\(c_1+N\)}{\(2 c_1\)}{ 1, 0}
                \reactionterm[\(k_{-2} c_2^2\)]{\(k_2 n c_2\)}{\(c_2+N\)}{\(2 c_2\)}{ 0, 1}
            \end{reactiontable}
        }{
            \begin{reactiontable}
                \reactionterm[\(w_2 c_2\)]{\(w_1 c_1\)}{\(c_1\)}{\(c_2\)}{-1, 1}
                \reactionterm[\(k_{-1} c_1^2\)]{\(k_1 n c_1\)}{\(c_1+N\)}{\(2 c_1\)}{ 1, 0}
                \reactionterm[\(k_{-2} c_2^2\)]{\(k_2 n c_2\)}{\(c_2+N\)}{\(2 c_2\)}{ 0, 1}
                \reactionterm{\(k_{1}c_1c_2\)}{\(c_1+c_2\)}{\(c_1+N\)}{-1, 0}
                \reactionterm{\(k_{2}c_1c_2\)}{\(c_1+c_2\)}{\(c_2+N\)}{0, -1}
            \end{reactiontable}
        }

        \column{0.5\textwidth}

        \only<1,2>{
            \begin{itemize}
                \item We can think about smothering them
            \end{itemize}
            \begin{align*}
                c_1 + c_2 &\xrightarrow{k_1} c_1 \\
                c_1 + c_2 &\xrightarrow{k_2} c_2
            \end{align*}
        }

    \end{columns}


    \begin{tikzpicture}[overlay, remember picture]

        \node<3-> at (frame cs:0.55, 0.45) {
                \includegraphics[width=0.38\framewidth, page=1]
                {ode/ssa:ode-5-2:I19:C2fp:T17475890580172482.pdf}
            };

        \node<3-> at (frame cs:0.55,-0.35) {
                \includegraphics[width=0.38\framewidth, page=3]
                {ode/ssa:ode-5-2:I19:C2fp:T17475890580172482.pdf}
            };

    \end{tikzpicture}

\end{frame}

\begin{frame}{Weird, what about...}

    \begin{columns}[T]
        \column{0.5\textwidth}

        \begin{itemize}
            \item With nutrition?
        \end{itemize}

        \alt<1>{
            \begin{reactiontable}
                \reactionterm[\(w_2 c_2\)]{\(w_1 c_1\)}{\(c_1\)}{\(c_2\)}{-1, 1}
                \reactionterm[\(k_{-1} c_1^2\)]{\(k_1 n c_1\)}{\(c_1+N\)}{\(2 c_1\)}{ 1, 0}
                \reactionterm[\(k_{-2} c_2^2\)]{\(k_2 n c_2\)}{\(c_2+N\)}{\(2 c_2\)}{ 0, 1}
                \reactionterm{\(k_{1}c_1c_2\)}{\(c_1+c_2\)}{\(c_1+N\)}{-1, 0}
                \reactionterm{\(k_{2}c_1c_2\)}{\(c_1+c_2\)}{\(c_2+N\)}{0, -1}
            \end{reactiontable}
        }{
            \begin{reactiontable}
                \reactionterm[\(w_2 c_2\)]{\(w_1 c_1\)}{\(c_1\)}{\(c_2\)}{-1, 1, 0}
                \reactionterm[\(k_{1} c_1^2\)]{\(k_1 n c_1\)}{\(c_1+N\)}{\(2 c_1\)}{ 1, 0, -1}
                \reactionterm[\(k_{2} c_2^2\)]{\(k_2 n c_2\)}{\(c_2+N\)}{\(2 c_2\)}{ 0, 1, -1}
                \reactionterm{\(k_{1}c_1c_2\)}{\(c_1+c_2\)}{\(c_1+N\)}{-1, 0, 1}
                \reactionterm{\(k_{2}c_1c_2\)}{\(c_1+c_2\)}{\(c_2+N\)}{0, -1, 1}
            \end{reactiontable}
        }

        \column{0.5\textwidth}


    \end{columns}


    \begin{tikzpicture}[overlay, remember picture]

        \node<2-> at (frame cs:0.55, 0.45) {
                \includegraphics[width=0.38\framewidth, page=1]
                {ode/ssa:ode-5-3:I19:C2:T1747490399601534.pdf}
            };

        \node<2-> at (frame cs:0.55,-0.35) {
                \includegraphics[width=0.38\framewidth, page=3]
                {ode/ssa:ode-5-3:I19:C2:T1747490399601534.pdf}
            };

    \end{tikzpicture}


\end{frame}


\subsection{My Proposed Model}

\begin{frame}{Two Species Compartment Model}

\end{frame}


\begin{frame}{My Model}

    \raggedbottom


    \begin{itemize}
        \item Built from principles
    \end{itemize}

    \tikz[overlay, remember picture] \node<2> at (frame cs:0,0)
    {%
        \begin{tikzpicture}
            \pgfmathsetnewmacro\genelength{35}
            \pgfmathsetnewmacro\geneheight{2}

            \pgfmathsetnewmacro\translength{10}
            \pgfmathsetnewmacro\transheight{2}

            \pgfmathsetnewmacro\reactionthickness{1}
            \pgfmathsetnewmacro\dnawidth{4}

            \pgfmathsetnewmacro\shortendist{6}
            \pgfmathsetnewmacro\bendstrength{1.05} % √(1²+(0.2-0.75)²)

            \pgfmathsetnewmacro\geneshift{5}

            \providecommand*\genecolor{red!80!black}
            \providecommand*\transcolor{green!80!black}

            \useasboundingbox (-4,-1.25) rectangle (5.5,1.35);
            \coordinate (O) at (0,0);

            \node[circle] (west) at (-3.5, 0) {};
            \node[circle] (east) at ( 4.5, 0) {};

            \node[anchor=south east] at (west) {\(5'\)};
            \node[anchor=south west] at (east) {\(3'\)};

            \draw[line width = \dnawidth] (west.center) to (east.center);

            \draw[line width = \dnawidth] (east.north) to (east.south);
            \draw[line width = \dnawidth] (west.north) to (west.south);


            \coordinate (g1) at (-2,0); \coordinate (g2) at ( 0,0);
            \coordinate (g3) at ( 3,0);

            \coordinate (t1) at (-3,0); \coordinate (t2) at (-1,0);
            \coordinate (t3) at ( 1,0);
            \coordinate (t4) at ( 2,0); \coordinate (t5) at ( 4,0);

            \draw[line width = \dnawidth+1, BackgroundColor, dashed] (t3) -- (t4);
            \draw[line width = \dnawidth+1, BackgroundColor] ([xshift=0.37cm]t3) to ([xshift=-0.43cm]t4);

            \foreach \POINT in {1,2,3}{
                \node[rectangle, fill=\genecolor, draw=\genecolor, minimum width=\genelength, minimum height=\geneheight] at (g\POINT) {};
            }

            \foreach \POINT in {1,2,3,4,5}{
                \node[rectangle, fill=\transcolor, draw=\transcolor, minimum width=\translength, minimum height=\transheight] at (t\POINT) {};
            }

            \node[above, yshift=5] (x1) at (g1) {\(X\)};
            \node[above, yshift=5] (x2) at (g2) {\(X\)};
            \node[above, yshift=5] (x3) at (g3) {\(X\)};

            \node  (B) at ( 4.75, 0.85) {\(B\)};
            \node (B2) at ( 0.50, 0.85) {\(B\)};

            \node  (N) at ( 2.00,-0.75) {\(N\)};
            \node (N2) at (-2.00, 1.00) {\(N\)};
            \node (N3) at ( 2.75, 1.00) {\(N\)};

            \node  (y) at ( 4.00,-0.75) {\(Y\)};
            \node (y2) at ( 1.50, 0.30) {\(Y\)};
            \node (y3) at ( 0.75,-0.75) {\(Y\)};

            \begin{scope}[line width=\reactionthickness, shorten <=\shortendist, shorten >=\shortendist]
                \draw[->] ([yshift=\geneshift]g1)
                    to[bend left] node[midway, sloped] (k1path) {}
                    ([yshift=\geneshift]g2);
                \draw[->] ([yshift=-\geneshift]g2)
                    to[bend left] node[midway] (k-1path) {}
                    ([yshift=-\geneshift]g1);
            \end{scope}

            \node[above] at (k1path) {\(k_1\)};
            \node[below] at (k-1path) {\(k_{-1}\)};

            \draw[line width=\reactionthickness, shorten >=6pt]
                (N2) to[bend right] ([yshift=2.3]k1path.center);

            % \draw[->,line width=\reactionthickness, shorten <=5.5pt]
            %     ([yshift=-2.6]k-1path.center) to[bend right] (N3);



            \begin{scope}[line width = \reactionthickness]
                \draw[shorten >=5pt, shorten <=4pt ]
                (B.center) -- (x3.center)
                node[pos=0.5, sloped] (midbx) {\(\bigtimes\)};
                \draw[shorten >=5pt, shorten <=4pt ]
                (B2.center) -- (y2.center)
                node[pos=0.5, sloped] (midby) {\(\bigtimes\)};

                \draw[->] (midbx.center) to[in=0, out=105, looseness=1.5] (N3);
                \draw[->] (midbx.center) to[in=160, out=105, looseness=3] (B);

                \draw[->] (midby.center) to[in=180, out=71, looseness=1.5] (N3);

                % \draw[->] (midxb.center) -- ++(105:0.5);
                % \draw[->] (midxb.center) to[bend right] (N2.east);

                \draw[->] (N) .. controls ++( 45:\bendstrength) and ++(135:\bendstrength) .. (y)
                node[midway, below] {\(k_2\)};
                \draw[->] (y3) -- (N) node[midway, above] {\(k_5\)};
            \end{scope}

            \node[anchor=north west] at (midbx) {\(k_3\)};
            \node[anchor=west] at ([yshift=2,xshift=1]midby) {\(k_4\)};


        \end{tikzpicture}
    };


\end{frame}


\section{Discussion}

\begin{frame}{ODE to SDE Differences}
\end{frame}



\section{End}
\begin{frame}{Questions?}
    \large Questions?
\end{frame}


\section{Bibliography}
\begin{frame}{Bibliography}

    \printbibliography

\end{frame}

\end{document}
